\SourceLeaf{271}%
Anwendung auf eine gewisse Zone zu betrachten sind, gern mit besonderer Sorgfalt und Schärfe auszuführen wünschen, und es verdienen daher einige dazu dienliche Umformungen hier einen Platz: eine Umformung wird ohnehin \textit{nothwendig}, wenn man von einer bestimmten Normalbreite nicht auf dem Ellipsoid sondern auf der Kugel, also von einem gegebenen Werthe von \(Q\) ausgehen, und daraus die übrigen Grössen berechnen will.

Führt man drei Hülfswinkel \(\varphi,\zeta,\eta\) ein, so dass
\[
\sin\varphi=e \quad .\ .\ .\ .\ .\ .\ .\ .\ .\ .\ .\ .\ .\ .\ .\quad(13)
\]
\[
\operatorname{tang}\zeta=\operatorname{tang}\varphi\cos P^2 \quad .\ .\ .\ .\ .\ .\ .\ .\ .\ .\quad(14)
\]
\[
\operatorname{tang}\eta=\sin\zeta\operatorname{tang}P \quad .\ .\ .\ .\ .\ .\ .\ .\ .\ .\ .\quad(15)
\]
so wird
\[
\alpha=\frac{1}{\cos\zeta} \quad .\ .\ .\ .\ .\ .\ .\ .\ .\ .\ .\ .\ .\ .\quad(16)
\]
\[
\sin Q=\cos\zeta\sin P \quad .\ .\ .\ .\ .\ .\ .\ .\ .\ .\ .\ .\quad(17)
\]
\[
\cos\eta\cos Q=\cos P \quad .\ .\ .\ .\ .\ .\ .\ .\ .\ .\ .\ .\ .\quad(18)
\]
\[
\sin\eta=\operatorname{tang}\zeta\operatorname{tang}Q \quad .\ .\ .\ .\ .\ .\ .\ .\ .\ .\quad(19)
\]
\[
\operatorname{tang}\tfrac{1}{2}(P-Q)=\operatorname{tang}\tfrac{1}{2}\zeta\mathord{.}\operatorname{tang}\tfrac{1}{2}\eta \quad .\ .\ .\ .\ .\ .\ .\ .\ .\quad(20)
\]
\[
\sin(2\zeta-\varphi)=e\cos2Q \quad .\ .\ .\ .\ .\ .\ .\ .\ .\ .\ .\ .\quad(21)
\]

Die Gleichung 18 folgt leicht aus der Verbindung von 15 und 17; sodann 19 aus der Verbindung von 15, 17, 18; ferner 20 aus 17, 18, 19, endlich 21 aus 14 und 17.

Am schärfsten wird man rechnen, wenn man, in dem Falle wo \(P\) gegeben ist, sich der Formeln 14, 15, 20 bedient, um der Reihe nach \(\zeta,\eta,Q\) zu bestimmen; für den Fall hingegen, wo \(Q\) gegeben ist, vermittelst der Gleichungen 21, 19, 20 die Werthe von \(\zeta,\eta,P\) ableitet: zur Controlle mag man dann noch eine oder einige der übrigen Gleichungen benutzen. Führt man noch einen vierten Hülfswinkel \(\theta\) ein, nach der Formel
\[
\sin\theta=e\sin P \quad .\ .\ .\ .\ .\ .\ .\ .\ .\ .\ .\ .\ .\ .\quad(22)
\]
so wird
\[
\cos\varphi=\cos\zeta\cos\eta\cos\theta \quad .\ .\ .\ .\ .\ .\ .\ .\ .\ .\quad(23)
\]
und die Formeln 9 und 12 erhalten folgende Gestalt:
\[
A=\frac{a\cos P}{\alpha\cos\theta\cos Q}=\frac{a\cos\eta}{\alpha\cos\theta}=\frac{a\cos\varphi}{\cos\theta^2}=\frac{a\cos\varphi}{1-ee\sin P^2}
\]
\[
k=\frac{\operatorname{tang}(45^0+\frac{1}{2}P)^\alpha}{\operatorname{tang}(45^0+\frac{1}{2}Q)\operatorname{tang}(45^0+\frac{1}{2}\theta)^{\alpha e}}
\]
